\section{قوانین مقیاس پذیری (\lr{Scaling Laws})}

\subsection{مقدمه}
قوانین مقیاس‌پذیری (\lr{Scaling Laws}) به عنوان یکی از ارکان اصلی در توسعه و درک مدل‌های زبانی بزرگ (\lr{LLMs}) شناخته می‌شوند. این قوانین که در ابتدا توسط محققانی چون \lr{Kaplan} و \lr{Hoffmann} (قانون \lr{Chinchilla}) مطرح شدند، نشان می‌دهند که عملکرد مدل‌ها با افزایش پارامترها، حجم داده‌ها و بودجه محاسباتی به صورت قابل پیش‌بینی بهبود می‌یابد. با این حال، تحقیقات اخیر نشان داده‌اند که قوانین کلاسیک دارای محدودیت‌هایی هستند؛ از جمله نادیده گرفتن هزینه‌های استنتاج (\lr{Inference})، تاثیر شکل معماری (\lr{Architecture Shape})، کیفیت داده‌ها و ساختارهای نوینی مانند «ترکیبی از خبرگان» (\lr{Mixture-of-Experts} یا \lr{MoE}). این بخش، با ترکیب و تحلیل ۱۹ مقاله علمی جدید، به بررسی ابعاد پنهان و توسعه‌یافته‌ی قوانین مقیاس‌پذیری می‌پردازد.

\subsection{اصلاح و توسعه قوانین مقیاس‌پذیری کلاسیک}
قوانین اولیه مانند \lr{Chinchilla} فرض می‌کردند که نسبت بهینه داده به پارامتر ($D/N$) یک مقدار ثابت (حدود ۲۰) است. اما لی و همکاران \cite{li2025farseer} با معرفی قانون \lr{Farseer} نشان می‌دهند که این نسبت ثابت نیست و با افزایش بودجه محاسباتی، نسبت بهینه $D/N$ نیز افزایش می‌یابد. روش \lr{Farseer} با استفاده از برازش قطعه‌ای دیفرانسیلی (\lr{Differential Piecewise Fitting})، خطای برون‌یابی (\lr{Extrapolation Error}) را نسبت به قانون \lr{Chinchilla} تا ۴۳۳ درصد کاهش می‌دهد.

از سوی دیگر، چن و همکاران \cite{chen2025revisiting} پدیده \lr{Sub-scaling} (کاهش نرخ بهبود عملکرد) را بررسی می‌کنند. این پدیده زمانی رخ می‌دهد که مدل‌ها بیش از حد آموزش می‌بینند (\lr{Over-training}) یا چگالی داده‌ها (افزونگی اطلاعات) بسیار بالا است. این مقاله یک قانون مقیاس‌پذیری زیر-بهینه (\lr{Sub-optimal}) ارائه می‌دهد که نشان می‌دهد داده‌های با چگالی بالا منجر به رشد توانی (کاهش بازدهی نهایی) و داده‌های با چگالی پایین منجر به رشد خطی در عملکرد می‌شوند.

همچنین، مک‌لیش و همکاران \cite{mcleish2025gemstones} با معرفی مجموعه \lr{Gemstones} (شامل بیش از ۴۰۰۰ چک‌پوینت)، شکنندگی قوانین مقیاس‌پذیری نسبت به انتخاب‌های ابرپارامتری (مانند نرخ یادگیری و شکل مدل) را برجسته می‌کنند. آن‌ها روش \lr{Convex Hull} را برای برازش قوانین پیشنهاد می‌دهند که نسبت به روش‌های پیشین بسیار مقاوم‌تر است.

\subsection{تاثیر معماری و شکل مدل (\lr{Model Shape \& Architecture})}
یکی از مهم‌ترین نقایص قوانین گذشته، نادیده گرفتن تاثیر شکل مدل (نسبت عمق به عرض) و هزینه‌های استنتاج است.

\begin{itemize}
    \item \textbf{عرض در برابر عمق:} بیان و همکاران \cite{bian2025scaling} نشان می‌دهند که معماری مدل مستقیماً بر تاخیر استنتاج (\lr{Inference Latency}) تاثیر می‌گذارد. مدل‌های عریض‌تر و کم‌عمق‌تر (\lr{Wider and Shallower}) با بودجه پارامتری یکسان، توان عملياتی (\lr{Throughput}) بسیار بالاتری در استنتاج دارند. در همین راستا، لیو و همکاران \cite{liu2026inverse} پدیده \lr{Inverse Depth Scaling} را معرفی می‌کنند؛ جایی که خطا با نسبت معکوس عمق ($1/\ell$) کاهش می‌یابد. آن‌ها دلیل این امر را میانگین‌گیری گروهی (\lr{Ensemble Averaging}) لایه‌ها می‌دانند، به این معنا که لایه‌های متعدد در \lr{LLM}ها کارهای مشابهی انجام می‌دهند و صرفاً با میانگین‌گیری، خطا را کاهش می‌دهند که نشان‌دهنده استفاده ناکارآمد از عمق است.
    
    \item \textbf{تعامل عصبی (\lr{Neural Interaction}):} سان و همکاران \cite{sun2026law} با استفاده از مفهوم \lr{Average Gradient Outer Product (AGOP)}، قانون «تعامل عصبی» را مطرح می‌کنند. آن‌ها نشان می‌دهند که تعمیم‌پذیری خوب نیازمند برهم‌نهی خوش‌خیم (\lr{Benign Superposition}) است؛ جایی که مدل با انرژی تعاملی پایین، سهم بالایی از حساسیت تعاملی را به دست می‌آورد. نسبت عمق به عرض ($R_{D/W}$) ابزاری برای قرار دادن مدل در این بازه بهینه است.
    
    \item \textbf{قوانین مقیاس‌پذیری طیفی (\lr{Spectral Scaling Laws}):} جها و ریگن \cite{jha2025spectral} نشان می‌دهند که افزایش عرض شبکه‌های پیش‌خور (\lr{FFN}) عمدتاً ظرفیت‌های کم‌انرژی (دم طیف) را گسترش می‌دهد، در حالی که زیرفضاهای غالب به سرعت اشباع می‌شوند. این قانون مقیاس‌پذیری طیفی نامتقارن نشان می‌دهد که انتخاب عرض \lr{FFN} یک مبادله بین ظرفیت دم و ظرفیت حالت غالب است.
    
    \item \textbf{معماری‌های جایگزین (\lr{xLSTM} و \lr{Mamba}):} بک و همکاران \cite{beck2026xlstm} قوانین مقیاس‌پذیری معماری \lr{xLSTM} (با پیچیدگی زمانی خطی) را با ترانسفورمرها مقایسه می‌کنند. نتایج نشان می‌دهد که \lr{xLSTM} در مبادله محاسبات-خطا (\lr{Compute-Loss Trade-off}) بر ترانسفورمرها برتری پارتو (\lr{Pareto-dominate}) دارد. همچنین، در مقاله‌ای دیگر \cite{anonymous2026recall} با رویکرد تفسیرپذیری مکانیسمی، نشان داده شده است که معماری \lr{Mamba} بازیابی اطلاعات (\lr{Associative Recall}) را از طریق یادگیری ضمنی توابع هش خطی انجام می‌دهد و کران‌های بالایی برای اندازه واژگان و ابعاد حالت بر اساس لم \lr{Johnson-Lindenstrauss} ارائه می‌دهد.
\end{itemize}

\subsection{قوانین مقیاس‌پذیری برای مدل‌های ترکیبی از خبرگان (\lr{MoE})}
مدل‌های \lr{MoE} به دلیل جداسازی تعداد کل پارامترها از هزینه محاسباتی (\lr{FLOPs})، نیازمند قوانین مقیاس‌پذیری مختص به خود هستند.

\begin{itemize}
    \item \textbf{اهرم کارایی (\lr{Efficiency Leverage}):} تیان و همکاران \cite{tian2025towards} معیار \lr{EL} را برای سنجش برتری محاسباتی \lr{MoE} نسبت به مدل‌های متراکم (\lr{Dense}) معرفی می‌کنند. آن‌ها نشان می‌دهند که نسبت فعال‌سازی (\lr{Activation Ratio}) محرک اصلی کارایی است و دانه‌بندی خبرگان (\lr{Expert Granularity}) به عنوان یک تعدیل‌کننده غیرخطی عمل می‌کند.
    
    \item \textbf{مبادله پارامتر و \lr{FLOPs}:} آبنار و همکاران \cite{abnar2025parameters} نشان می‌دهند که برای یک بودجه محاسباتی ثابت، افزایش سطح پراکندگی (\lr{Sparsity}) — یعنی داشتن پارامترهای کل بیشتر اما پارامترهای فعال کمتر — منجر به کاهش خطای آموزش می‌شود.
    
    \item \textbf{کارایی حافظه:} برخلاف تصور رایج، لودزیجوسکی و همکاران \cite{ludziejewski2025joint} ثابت می‌کنند که مدل‌های \lr{MoE} می‌توانند از نظر حافظه نیز کارآمدتر از مدل‌های متراکم باشند، به شرطی که روی توکن‌های بیشتری آموزش ببینند.
    
    \item \textbf{بازیافت مدل‌ها (\lr{Upcycling}):} لیو و همکاران \cite{liew2025scaling} به بررسی تبدیل مدل‌های متراکم از پیش‌آموزش‌دیده به \lr{MoE} می‌پردازند. آن‌ها یک قانون مقیاس‌پذیری مشترک کشف کردند که نشان می‌دهد \lr{Upcycling} تنها زمانی کارآمدتر از آموزش از صفر (\lr{From-scratch}) است که هزینه اولیه (\lr{Sunk Cost} - توکن‌های آموزش مدل متراکم) و اندازه مدل نسبتاً کوچک باشند.
\end{itemize}

\subsection{نقش داده‌ها: کیفیت، ترکیب و چگالی}
داده‌ها سوخت اصلی قوانین مقیاس‌پذیری هستند و کیفیت آن‌ها به اندازه کمیتشان اهمیت دارد.

\begin{itemize}
    \item \textbf{کیفیت داده‌ها:} سوبرامانیام و همکاران \cite{subramanyam2026scaling} یک پارامتر بدون بعد کیفیت به نام $Q \in (0, 1]$ را معرفی کرده و قانون \lr{Chinchilla} را برای در نظر گرفتن کیفیت داده‌ها بسط می‌دهند. آن‌ها نشان می‌دهند که داده‌های با کیفیت بالا می‌توانند نیاز به اندازه مدل و محاسبات را به شدت کاهش دهند و اندازه موثر نمونه به صورت $D_{eff} = D \cdot Q^\gamma$ مقیاس می‌یابد.
    
    \item \textbf{ترکیب بهینه داده‌ها (\lr{Data Mixtures}):} انتخاب نسبت ترکیب دامنه‌های مختلف داده معمولاً با آزمون و خطا انجام می‌شود. شوکور و همکاران \cite{shukor2025scaling} روشی سیستماتیک ارائه می‌دهند که با استفاده از برازش قوانین مقیاس‌پذیری روی اجراهای مقیاس کوچک، می‌توان وزن‌های بهینه دامنه را برای آموزش مدل‌های بسیار بزرگ پیش‌بینی کرد.
    
    \item \textbf{ساختار سلسله‌مراتبی داده‌ها:} کاگنتا و همکاران \cite{cagnetta2025scaling} با استفاده از مدل سلسله‌مراتب تصادفی (\lr{RHM})، نشان می‌دهند که شبکه‌های عصبی پیچیدگی‌های داده را لایه به لایه یاد می‌گیرند. آن‌ها ثابت می‌کنند که شبکه‌های پیچشی (\lr{CNNs}) به دلیل اشتراک وزن و اتصال محلی، در یادگیری داده‌های سلسله‌مراتبی سریع‌تر از ترانسفورمرها مقیاس می‌یابند.
\end{itemize}

\subsection{فراتر از خطای مطلق: معیارهای نوین ارزیابی}
ارزیابی مدل‌ها صرفاً بر اساس خطای آنتروپی متقاطع (\lr{Cross-Entropy}) تصویر کاملی از توانایی‌های مدل ارائه نمی‌دهد.

\begin{itemize}
    \item \textbf{قوانین مقیاس‌پذیری نسبی (\lr{Relative Scaling Laws}):} هلد و همکاران \cite{held2026relative} استدلال می‌کنند که ارزیابی‌های تجمیعی، تفاوت عملکرد در زیرگروه‌ها را پنهان می‌کند. آن‌ها قوانین نسبی را معرفی می‌کنند تا نشان دهند چگونه شکاف عملکرد بین توزیع‌ها با افزایش مقیاس تغییر می‌کند. نتایج نشان می‌دهد که مقیاس‌پذیری یک تساوی‌بخش جهانی نیست؛ در حالی که دامنه‌های علمی در \lr{MMLU} به همگرایی می‌رسند، خطرات هوش مصنوعی (مانند توانایی فریبکاری) با افزایش مقیاس واگرا شده و افزایش می‌یابند.
    
    \item \textbf{احتمال مبتنی بر رتبه (\lr{Relative-Based Probability}):} یو و همکاران \cite{yue2025relative} معیار \lr{RBP} را معرفی می‌کنند که احتمال قرار گرفتن توکن صحیح در میان $k$ پیش‌بینی برتر را می‌سنجد. آن‌ها نشان می‌دهند که $-\log(RBP_k)$ نیز از یک قانون توانی پیروی می‌کند. این دیدگاه پدیده ظهور توانایی‌ها (\lr{Emergence}) را نه به عنوان یک جهش ناگهانی، بلکه به عنوان یک اثر ماکروسکوپی قابل پیش‌بینی از مقیاس‌پذیری میکروسکوپی هموار توضیح می‌دهد.
    
    \item \textbf{قانون چگال‌سازی (\lr{Densing Law}):} در نهایت، شیائو و همکاران \cite{xiao2025densing} برای ارزیابی روند پیشرفت کارایی \lr{LLM}ها در طول زمان، مفهوم چگالی قابلیت (\lr{Capability Density}) را معرفی می‌کنند. مشاهدات تجربی آن‌ها نشان می‌دهد که حداکثر چگالی قابلیت مدل‌های متن‌باز تقریباً هر ۳.۵ ماه دو برابر می‌شود که نشان‌دهنده کاهش نمایی هزینه‌های استنتاج و نیازهای پارامتری برای رسیدن به یک سطح عملکرد ثابت است.
\end{itemize}

\subsection{نتیجه‌گیری}
بررسی جامع این مقالات نشان‌دهنده یک تغییر پارادایم در درک ما از قوانین مقیاس‌پذیری است. ما از دوران «بزرگتر بودن همیشه بهتر است» (صرفاً با افزایش پارامتر و داده) عبور کرده‌ایم. قوانین مقیاس‌پذیری نوین اکنون چندوجهی شده‌اند و باید هزینه‌های استنتاج، شکل معماری (عمق/عرض)، کیفیت و چگالی داده‌ها، ساختارهای پراکنده (\lr{MoE}) و معیارهای ارزیابی نسبی را در بر بگیرند. این بینش‌های جدید به محققان و مهندسان اجازه می‌دهد تا مدل‌های زبانی آینده را نه تنها قدرتمندتر، بلکه از نظر محاسباتی، حافظه و استنتاج به مراتب کارآمدتر و هدفمندتر طراحی کنند.